% ---------------------------------------------------------------------------
% Framework chapter. Requires amsmath, amssymb, amsthm, booktabs, tikz.
% The universes \Utime, \Uact, \Uobj, \Utype, \Uattr, \Uval, \Uqual are
% assumed from the preliminaries; the macros below are only fallbacks.
% ---------------------------------------------------------------------------
\providecommand{\mi}[1]{\mathit{#1}}
\providecommand{\la}{\langle}
\providecommand{\ra}{\rangle}
\providecommand{\Utime}{\mathcal{U}_{\mi{time}}}
\providecommand{\Uact}{\mathcal{U}_{\mi{act}}}
\providecommand{\Uobj}{\mathcal{U}_{\mi{obj}}}
\providecommand{\Utype}{\mathcal{U}_{\mi{type}}}
\providecommand{\Uattr}{\mathcal{U}_{\mi{attr}}}
\providecommand{\Uval}{\mathcal{U}_{\mi{val}}}
\providecommand{\Uqual}{\mathcal{U}_{\mi{qual}}}
\providecommand{\Uev}{\mathcal{U}_{\mi{ev}}}
\providecommand{\Uoo}{\mathcal{U}_{\mi{o2o}}}
\providecommand{\Uau}{\mathcal{U}_{\mi{au}}}
\providecommand{\Uitem}{\mathcal{U}_{\mi{item}}}
\providecommand{\Umod}{\mathcal{U}_{\mi{model}}}
\providecommand{\Uhist}{\mathcal{U}_{\mi{hist}}}
\providecommand{\pfun}{\nrightarrow}
\providecommand{\restr}{\mathord{\upharpoonright}}

\section{A Buffer-Based Streaming Framework}
\label{sec:framework}

This section introduces a framework for object-centric streaming algorithms
with bounded memory.
The input is an object-centric event stream.
The output is a process model after every stream item.
Everything an algorithm remembers is stored in \emph{buffers}.
Each buffer has a fixed capacity.
When a buffer overflows, a \emph{heuristic} selects the key to evict.

The framework separates two concerns.
An \emph{algorithm} defines what is stored and how the model is computed.
A \emph{configuration} defines how much is stored and what is evicted.
The same algorithm can therefore be run with different capacities and
heuristics, e.g., the cache policies of Liss et al.~\cite{Liss2025} or the
lifecycle estimation of M\"uhlmeyer et al.~\cite{Muehlmeyer2026}.
Figure~\ref{fig:framework} shows how one stream item is processed.

\begin{figure}[t]
\centering
\begin{tikzpicture}[
    node distance=6mm and 6mm,
    box/.style={draw, rounded corners, minimum height=8mm, align=center, font=\small},
    lbl/.style={font=\footnotesize, align=center},
    >=stealth]
  \node[lbl] (item) {item\\$s_i$};
  \node[box, right=of item] (upd) {$\mi{upd}$};
  \node[box, right=of upd] (ev1) {evict $B_1$\\[-1pt]{\footnotesize $(c_1, h_1)$}};
  \node[lbl, right=of ev1] (dots) {$\cdots$};
  \node[box, right=of dots] (evn) {evict $B_n$\\[-1pt]{\footnotesize $(c_n, h_n)$}};
  \node[lbl, right=of evn] (state) {state\\$q_i$};
  \node[box, below=of state] (model) {$\mi{model}$};
  \node[lbl, left=of model] (out) {model after $s_i$};
  \node[lbl, above=7mm of upd] (prev) {state $q_{i-1}$};
  \draw[->] (item) -- (upd);
  \draw[->] (prev) -- (upd);
  \draw[->] (upd) -- (ev1);
  \draw[->] (ev1) -- (dots);
  \draw[->] (dots) -- (evn);
  \draw[->] (evn) -- (state);
  \draw[->] (state) -- (model);
  \draw[->] (model) -- (out);
\end{tikzpicture}
\caption{Processing of stream item $s_i$.
The algorithm first updates its state with $\mi{upd}$.
Then the buffers $B_1, \dots, B_n$ are handled in order.
While buffer $B_j$ holds more than $c_j$ keys, heuristic $h_j$ selects a key
and $\mi{fin}$ evicts it.
Evicting a key may record information in later buffers.
The resulting state $q_i$ yields the current model.
The algorithm provides $\mi{upd}$, $\mi{fin}$, and $\mi{model}$; the
configuration provides the capacities $c_j$ and heuristics $h_j$.}
\label{fig:framework}
\end{figure}

\subsection{Object-Centric Event Streams}
\label{sec:framework:streams}

We use the following notation.
$X \pfun Y$ is the set of partial functions from $X$ to $Y$.
For $f \in X \pfun Y$, $f \oplus g$ updates $f$ with the pairs in $g$, and
$f \restr X'$ restricts $f$ to $\mi{dom}(f) \cap X'$.
For $b \in K \pfun \mathbb{N}$ and $k \in K$, we write $b[k] = b(k)$ if
$k \in \mi{dom}(b)$ and $b[k] = 0$ otherwise.
For a finite multiset $M$ over $K$, $b \uplus M$ adds the multiplicities of
$M$ to $b$, i.e., $\mi{dom}(b \uplus M) = \mi{dom}(b) \cup \{k \in M\}$ and
$(b \uplus M)(k) = b[k] + M(k)$.

We adapt the object-centric event streams of Liss et al.~\cite{Liss2025}.
In contrast to~\cite{Liss2025}, each position of the stream holds a single
item.
Simultaneous items follow each other with equal timestamps.

\begin{definition}[Stream Item]
\label{def:item}
$\Uev = \Uact \times \mathcal{P}(\Uobj \times \Utype \times \Uqual) \times
(\Uattr \pfun \Uval)$ is the universe of events,
$\Uoo = \Uobj \times \Utype \times \Uqual \times \Uobj \times \Utype$ is the
universe of object-to-object (O2O) updates, and
$\Uau = \Uobj \times \Utype \times \Uattr \times \Uval$ is the universe of
attribute updates.
$\Uitem = (\Uev \cup \Uoo \cup \Uau) \times \Utime$ is the universe of
stream items.
An event $(a, R, d) \in \Uev$ has activity $a$, a finite set $R$ of
qualified event-to-object relations, and attribute values $d$.

Let $s = (x, t) \in \Uitem$.
The typed objects of $s$ are
\begin{itemize}
  \item $\mi{tobj}(s) = \{(o, \mi{ot}) \mid \exists_{q \in \Uqual}\
        (o, \mi{ot}, q) \in R\}$ if $x = (a, R, d) \in \Uev$,
  \item $\mi{tobj}(s) = \{(o_1, \mi{ot}_1), (o_2, \mi{ot}_2)\}$ if
        $x = (o_1, \mi{ot}_1, q, o_2, \mi{ot}_2) \in \Uoo$, and
  \item $\mi{tobj}(s) = \{(o, \mi{ot})\}$ if
        $x = (o, \mi{ot}, \mi{attr}, v) \in \Uau$.
\end{itemize}
The objects of $s$ are
$\mi{obj}(s) = \{o \mid \exists_{\mi{ot} \in \Utype}\ (o, \mi{ot}) \in
\mi{tobj}(s)\}$.
\end{definition}

\begin{definition}[Object-Centric Event Stream]
\label{def:stream}
An object-centric event stream is a finite or infinite sequence
$S = \la s_1, s_2, \dots \ra$ of stream items $s_i = (x_i, t_i) \in \Uitem$.
For $i \ge 1$, $S_{\le i} = \la s_1, \dots, s_i \ra$ is the prefix of $S$
up to $s_i$.
$S$ is well-formed if and only if
\begin{enumerate}
  \item[(i)] $t_i \le t_j$ for all $1 \le i \le j$, and
  \item[(ii)] $\mi{ot} = \mi{ot}'$ for all $i, j \ge 1$, $o \in \Uobj$, and
        $\mi{ot}, \mi{ot}' \in \Utype$ with $(o, \mi{ot}) \in \mi{tobj}(s_i)$
        and $(o, \mi{ot}') \in \mi{tobj}(s_j)$.
\end{enumerate}
\end{definition}

A prefix $S_{\le i}$ is the event data that an offline algorithm would
receive after $s_i$, i.e., the materialized event log.

\begin{example}[Running Example]
\label{ex:running}
Table~\ref{tab:log} shows an event log of an order-management process with
the object types $\mi{order}$ and $\mi{item}$.
An order is placed together with its items.
Each item is picked, and the order is paid.
Finally, the order is shipped together with its items.
Orders $o_1$ and $o_2$ complete this lifecycle, while order $o_3$ and item
$i_3$ are still in progress after the last event.
Each event $e_i$ becomes the stream item $s_i = ((a_i, R_i, \emptyset), i)$
with timestamp $i$, where $R_i$ relates the objects of $e_i$ with the
qualifier $\mi{ref} \in \Uqual$, e.g.,
\[
  s_1 = \bigl((\mi{place}, \{(o_1, \mi{order}, \mi{ref}),
  (i_1, \mi{item}, \mi{ref})\}, \emptyset),\ 1\bigr).
\]
The resulting stream $S_{\mi{ex}} = \la s_1, \dots, s_9 \ra$ is well-formed.
\end{example}

\begin{table}[t]
\centering
\caption{Event log of the running example.
Event $e_i$ becomes stream item $s_i$ of the stream $S_{\mi{ex}}$.
Activities are abbreviated as $\mi{place}$ (place order), $\mi{pick}$
(pick item), $\mi{pay}$ (pay order), and $\mi{ship}$ (ship order).}
\label{tab:log}
\begin{tabular}{@{}clll@{}}
\toprule
Event & Activity & Orders & Items \\
\midrule
$e_1$ & $\mi{place}$ & $o_1$ & $i_1$ \\
$e_2$ & $\mi{place}$ & $o_2$ & $i_2$ \\
$e_3$ & $\mi{pick}$ & & $i_1$ \\
$e_4$ & $\mi{pay}$ & $o_1$ & \\
$e_5$ & $\mi{ship}$ & $o_1$ & $i_1$ \\
$e_6$ & $\mi{pick}$ & & $i_2$ \\
$e_7$ & $\mi{place}$ & $o_3$ & $i_3$ \\
$e_8$ & $\mi{pay}$ & $o_2$ & \\
$e_9$ & $\mi{ship}$ & $o_2$ & $i_2$ \\
\bottomrule
\end{tabular}
\end{table}

\subsection{Streaming Algorithms}
\label{sec:framework:algorithms}

\begin{definition}[Buffer, State]
\label{def:buffer}
A buffer $B = (K, V)$ consists of a set of keys $K$ and a set of values $V$.
Its content is a partial function $b \in K \pfun V$ with a finite domain.
Let $B_1, \dots, B_n$ be buffers with $B_j = (K_j, V_j)$.
$Q = (K_1 \pfun V_1) \times \dots \times (K_n \pfun V_n)$ is the set of
states, restricted to contents with finite domains, and
$q_\emptyset = (\emptyset, \dots, \emptyset)$ is the empty state.
For $q = (b_1, \dots, b_n) \in Q$, $1 \le j \le n$, $k \in K_j$, and
$v \in V_j$, $\mi{del}_j(q, k)$ is $q$ with $k$ removed from $b_j$, and
$\mi{put}_j(q, k, v)$ is $q$ with $b_j(k) = v$.
\end{definition}

\begin{definition}[Streaming Algorithm]
\label{def:algorithm}
Let $\Umod$ be a universe of models.
A streaming algorithm is a tuple
$A = (B_1, \dots, B_n, \mi{upd}, \mi{fin}, \mi{model})$ where
$B_1, \dots, B_n$ are buffers with states $Q$ and
\begin{itemize}
  \item $\mi{upd} \in Q \times \Uitem \rightarrow Q$ processes a stream item,
  \item $\mi{fin}(q, j, k) \in Q$ evicts key $k \in \mi{dom}(b_j)$ from
        buffer $B_j$ in state $q = (b_1, \dots, b_n)$, and
  \item $\mi{model} \in Q \rightarrow \Umod$ yields the current model,
\end{itemize}
such that $\mi{fin}(q, j, k)$ and $\mi{del}_j(q, k)$ coincide on the buffers
$B_1, \dots, B_j$.
\end{definition}

Evicting $k$ thus removes $k$ and leaves all earlier buffers unchanged.
It may change later buffers, e.g., to record the last activity of an evicted
object as an end activity.

\begin{definition}[Used Keys]
\label{def:used}
Let $A = (B_1, \dots, B_n, \mi{upd}, \mi{fin}, \mi{model})$ be a streaming
algorithm with states $Q$, let $q \in Q$, and let $s \in \Uitem$.
The keys of $B_j$ used by $s$ in $q$ are
\[
  \mi{used}_j(q, s) = \{k \in K_j \mid k \in \mi{dom}(b_j')
  \text{ with } \mi{upd}(\mi{del}_j(q, k), s) = (b_1', \dots, b_n')\}.
\]
\end{definition}

An item uses a key if it would create the key if it were missing.
For a missing key, being used means being created.
For the algorithms in this paper, the used keys are exactly the keys that the
item inserts, changes, or reads, e.g., the last activity of an object.
Condition~(i) of Proposition~\ref{prop:correctness} requires this in general.
We assume that $\mi{used}_j(q, s)$ is finite.

\begin{example}[Streaming Directly-Follows Algorithm]
\label{ex:dfg}
The algorithm $A_{\mi{df}}$ discovers per object type the directly-follows
arcs and the start and end activities.
It is the directly-follows part of~\cite{Muehlmeyer2026}.
It uses four buffers:
\begin{itemize}
  \item $B_1 = (\Uobj, \Utype \times \Uact)$ maps each object to its type and
        last activity,
  \item $B_2 = (\Uact \times \Uact \times \Utype, \mathbb{N})$ counts arcs,
  \item $B_3 = (\Utype \times \Uact, \mathbb{N})$ counts start activities, and
  \item $B_4 = (\Utype \times \Uact, \mathbb{N})$ counts the end activities
        of evicted objects.
\end{itemize}
Let $q = (b_1, b_2, b_3, b_4) \in Q$ and let $s = ((a, R, d), t)$ be an event
item with $T = \mi{tobj}(s)$.
Then $\mi{upd}(q, s) = (b_1', b_2', b_3', b_4)$ with
\begin{align*}
  b_1' &= b_1 \oplus \{(o, (\mi{ot}, a)) \mid (o, \mi{ot}) \in T\},\\
  b_2' &= b_2 \uplus [(a', a, \mi{ot}) \mid (o, \mi{ot}) \in T \ \wedge\
          b_1(o) = (\mi{ot}, a')], \text{ and}\\
  b_3' &= b_3 \uplus [(\mi{ot}, a) \mid (o, \mi{ot}) \in T \ \wedge\
          o \notin \mi{dom}(b_1)].
\end{align*}
Other items are ignored, i.e., $\mi{upd}(q, s) = q$.
Evicting an object records its last activity as an end activity:
\[
  \mi{fin}(q, 1, o) = \bigl(b_1 \restr (\mi{dom}(b_1) \setminus \{o\}),\
  b_2,\ b_3,\ b_4 \uplus [b_1(o)]\bigr).
\]
For $j \in \{2, 3, 4\}$, $\mi{fin}(q, j, k) = \mi{del}_j(q, k)$.
The model consists of the arcs, start activities, and end activities:
\[
  \mi{model}(q) = \bigl(\mi{dom}(b_2),\ \mi{dom}(b_3),\
  \mi{dom}(b_4) \cup \mi{rng}(b_1)\bigr).
\]
The end activities of objects that are still buffered are read from $b_1$.
For the running example, the state after $s_3$ without evictions is
$q_3 = (b_1, b_2, b_3, \emptyset)$ with
\begin{align*}
  b_1 &= \{(o_1, (\mi{order}, \mi{place})), (o_2, (\mi{order}, \mi{place})),
         (i_1, (\mi{item}, \mi{pick})), (i_2, (\mi{item}, \mi{place}))\},\\
  b_2 &= \{((\mi{place}, \mi{pick}, \mi{item}), 1)\}, \text{ and }
  b_3 = \{((\mi{order}, \mi{place}), 2), ((\mi{item}, \mi{place}), 2)\}.
\end{align*}
The stream item $s_4$ records the payment of $o_1$.
It uses $\mi{used}_1(q_3, s_4) = \{o_1\}$ and
$\mi{used}_2(q_3, s_4) = \{(\mi{place}, \mi{pay}, \mi{order})\}$, and no
keys of $B_3$ and $B_4$.
\end{example}

\subsection{Bounded Runs}
\label{sec:framework:runs}

\begin{definition}[Heuristic, Configuration]
\label{def:configuration}
Let $A = (B_1, \dots, B_n, \mi{upd}, \mi{fin}, \mi{model})$ be a
streaming algorithm with states $Q$.
$\Uhist = (\Uitem \times Q)^*$ is the set of histories.
A heuristic for buffer $B_j$ is a function $h \in K_j \times \Uhist
\rightarrow \mathbb{R}$.
The score $h(k, \eta)$ rates key $k$ in the last state of $\eta$.
A higher score means that $k$ is evicted sooner.
A configuration $C = ((c_1, h_1), \dots, (c_n, h_n))$ assigns to each buffer
$B_j$ a capacity $c_j \in \mathbb{N} \cup \{\infty\}$ and a heuristic $h_j$.
$C$ is unbounded if $c_j = \infty$ for all $1 \le j \le n$.
$C$ fits a stream $S = \la s_1, s_2, \dots \ra$ if and only if
$c_j \ge |\mi{used}_j(q, s_i)|$ for all $1 \le j \le n$, $q \in Q$, and
$i \ge 1$.
\end{definition}

A heuristic may use, for each key $k$ of $B_j$, its insertion position
$\mi{ins}_j^\eta(k)$, the position of its last use $\mi{last}_j^\eta(k)$, and
its number of uses since insertion $\mi{freq}_j^\eta(k)$.
These are derived from the history $\eta$
(Definition~\ref{def:bookkeeping} in Appendix~\ref{app:framework}).
A fitting configuration ensures that every single item of the stream fits
into every buffer.
For buffer $B_1$ of Example~\ref{ex:dfg}, the capacity must be at least the
largest number of objects in an event of the stream.

\begin{definition}[Run]
\label{def:run}
Let $S = \la s_1, s_2, \dots \ra$ be an object-centric event stream, let
$A = (B_1, \dots, B_n, \mi{upd}, \mi{fin}, \mi{model})$ be a
streaming algorithm with states $Q$, and let
$C = ((c_1, h_1), \dots, (c_n, h_n))$ be a configuration that fits $S$.
The run of $A$ on $S$ under $C$ is the sequence of states
$\la q_0, q_1, \dots \ra$ with $q_0 = q_\emptyset$ and, for $i \ge 1$,
\[
  q_i = \mi{ev}_n^i(\dots \mi{ev}_1^i(\mi{upd}(q_{i-1}, s_i)) \dots).
\]
Let $\eta = \la (s_1, q_1), \dots, (s_{i-1}, q_{i-1}) \ra$.
For $1 \le j \le n$ and $q = (b_1, \dots, b_n) \in Q$, let
$\eta_q = \eta \cdot \la (s_i, q) \ra$ and
\[
  \mi{ev}_j^i(q) =
  \begin{cases}
    q & \text{if } |\mi{dom}(b_j)| \le c_j,\\
    \mi{ev}_j^i(\mi{fin}(q, j, k^*)) & \text{otherwise,}
  \end{cases}
\]
where $k^* \in \mi{dom}(b_j) \setminus \mi{used}_j(q_{i-1}, s_i)$ maximizes
$h_j(k^*, \eta_q)$.
Ties are broken by the smallest $\mi{ins}_j^{\eta_q}(k^*)$ and then by a fixed
total order on $K_j$.
The model after $s_i$ is $\mi{model}(q_i)$.
The unbounded run of $A$ on $S$ is its run under an unbounded configuration.
\end{definition}

Each item is processed in two phases.
First, $\mi{upd}$ processes the item.
Then, each buffer evicts keys until it respects its capacity.
Keys used by the current item are never evicted.
Every eviction removes a key from $B_j$, and later evictions never insert into
$B_j$.
Hence, $B_j$ holds at most $c_j$ keys after every item.
The unbounded run never evicts.

Table~\ref{tab:heuristics} lists the cache replacement strategies and priority
policies of Liss et al.~\cite{Liss2025} as heuristics.
The lifecycle estimation of M\"uhlmeyer et al.~\cite{Muehlmeyer2026} evicts
an object once its inactivity exceeds its expected remaining lifetime.
We turn it into a heuristic that ranks objects by how far they exceed this
lifetime.

\begin{table}[t]
\centering
\caption{Cache replacement strategies and priority policies
of~\cite{Liss2025} as heuristics for buffer $B_j$.
Let $\eta \in \Uhist$ with $i = |\eta|$.
For a function $f$ over the keys of $B_j$, $\hat{f}$ denotes its min-max
normalization over the keys of $B_j$ in the last state of $\eta$, with
$\hat{f} = 0$ if all keys have the same value.}
\label{tab:heuristics}
\begin{tabular}{@{}ll@{}}
\toprule
Strategy & Score $h(k, \eta)$ \\
\midrule
First in, first out (FIFO) & $i - \mi{ins}_j^\eta(k)$ \\
Least recently used (LRU) & $i - \mi{last}_j^\eta(k)$ \\
Least frequently used (LFU) & $-\mi{freq}_j^\eta(k)$ \\
LFU with dynamic aging (LFU-DA) & $\hat{h}_{\mi{FIFO}}(k, \eta) +
  \hat{h}_{\mi{LFU}}(k, \eta)$ \\
Random replacement (RR) & independent and uniform in $[0, 1]$ \\
Priority policy $p$ with weight $w \in [0, 1]$ &
  $w \cdot \hat{h}(k, \eta) + (1 - w) \cdot \hat{p}(k, \eta)$ \\
\bottomrule
\end{tabular}
\end{table}

\begin{definition}[Lifecycle Heuristic]
\label{def:lifecycle}
Let $A_{\mi{df}}$ be the algorithm of Example~\ref{ex:dfg}, let
$T_{\mi{int}}, T_{\max} \in \mathbb{N}$ with $T_{\mi{int}} < T_{\max}$, and
let $z > 0$.
Let $\eta \in \Uhist$ with $i = |\eta|$ and last state
$(b_1, b_2, b_3, b_4)$.
For $o \in \mi{dom}(b_1)$ with $b_1(o) = (\mi{ot}, a)$, let
$n = \sum_{a' \in \Uact} b_4[(\mi{ot}, a')]$ be the number of completed
lifecycles of type $\mi{ot}$, let $n_a = b_4[(\mi{ot}, a)]$ be the number
of those ending with $a$, and let
\[
  p = \frac{n_a + 1}{n + 2}, \qquad
  v = \frac{(n_a + 1)(n - n_a + 1)}{(n + 2)^2 (n + 3)}, \qquad
  \mi{conf}_\eta(o) = \max\bigl(0,\ p - z \sqrt{v}\bigr).
\]
The expected lifetime of $o$ is
\[
  \mi{life}_\eta(o) =
  \begin{cases}
    T_{\max} & \text{if } \mi{conf}_\eta(o) = 0,\\
    T_{\mi{int}} + 1 & \text{if } \mi{conf}_\eta(o) \ge 0.9999,\\
    T_{\mi{int}} + \bigl\lceil (T_{\max} - T_{\mi{int}}) \cdot
      0.95^{\mi{conf}_\eta(o) / (1 - \mi{conf}_\eta(o))} \bigr\rceil
      & \text{otherwise.}
  \end{cases}
\]
The lifecycle heuristic for $B_1$ is
\[
  h_{\mi{LC}}(o, \eta) = \frac{i - \mi{last}_1^\eta(o)}{\mi{life}_\eta(o)}.
\]
\end{definition}

The confidence $\mi{conf}_\eta(o)$ is the lower bound of the Beta-binomial
credible interval used in~\cite{Muehlmeyer2026}.
It estimates how likely the last activity of $o$ ends a lifecycle of its type.
A score above~$1$ means that $o$ has been inactive for longer than expected.
In contrast to~\cite{Muehlmeyer2026}, the thresholds do not trigger
evictions, only the capacity does.
As long as $b_4$ is empty, all objects have the same expected lifetime, and
$h_{\mi{LC}}$ ranks objects like LRU.
Since $b_4$ is filled by evictions, the heuristic learns end activities only
once $B_1$ is full.

\begin{example}[Bounded Run]
\label{ex:run}
We run $A_{\mi{df}}$ on the stream $S_{\mi{ex}}$ of
Example~\ref{ex:running}.
The configuration uses $c_1 = 4$ and $h_1 = h_{\mi{LRU}}$ for $B_1$ and
$c_2 = c_3 = c_4 = \infty$.
Ties are broken by the order $o_1, o_2, o_3, i_1, i_2, i_3$.
Table~\ref{tab:run} shows the run.

\begin{table}[t]
\centering
\caption{Run of $A_{\mi{df}}$ on $S_{\mi{ex}}$ under $c_1 = 4$ and LRU.
Column $\mi{dom}(b_1)$ lists the buffered objects after each item.
The last column lists the evicted objects with the end activity recorded in
$b_4$.
The eviction of $o_2$ at $i = 7$ is premature, because $o_2$ occurs again at
$i = 8$.}
\label{tab:run}
\begin{tabular}{@{}clll@{}}
\toprule
$i$ & $s_i$ & $\mi{dom}(b_1)$ after $s_i$ & evicted (end activity) \\
\midrule
1 & $(\mi{place}, \{o_1, i_1\})$ & $o_1, i_1$ & -- \\
2 & $(\mi{place}, \{o_2, i_2\})$ & $o_1, o_2, i_1, i_2$ & -- \\
3 & $(\mi{pick}, \{i_1\})$ & $o_1, o_2, i_1, i_2$ & -- \\
4 & $(\mi{pay}, \{o_1\})$ & $o_1, o_2, i_1, i_2$ & -- \\
5 & $(\mi{ship}, \{o_1, i_1\})$ & $o_1, o_2, i_1, i_2$ & -- \\
6 & $(\mi{pick}, \{i_2\})$ & $o_1, o_2, i_1, i_2$ & -- \\
7 & $(\mi{place}, \{o_3, i_3\})$ & $o_3, i_1, i_2, i_3$ &
  $o_2$ ($\mi{place}$), $o_1$ ($\mi{ship}$) \\
8 & $(\mi{pay}, \{o_2\})$ & $o_2, o_3, i_2, i_3$ & $i_1$ ($\mi{ship}$) \\
9 & $(\mi{ship}, \{o_2, i_2\})$ & $o_2, o_3, i_2, i_3$ & -- \\
\bottomrule
\end{tabular}
\end{table}

The buffer first overflows at $i = 7$, when $o_3$ and $i_3$ arrive.
The candidates are $o_1$, $o_2$, $i_1$, and $i_2$, last used at the positions
$5$, $2$, $5$, and $6$.
LRU first evicts $o_2$ and then, by the tie-breaking order, $o_1$.
At $i = 8$, $o_2$ is created again as a new object, and LRU evicts $i_1$.
Table~\ref{tab:models} compares the models after $s_9$.
The bounded model contains the false start activity $\mi{pay}$ for orders.
The premature eviction of $o_2$ causes two further errors that are masked.
The arc $(\mi{place}, \mi{pay}, \mi{order})$ of $o_2$ is missing, but $o_1$
contributes it as well.
The false end activity $\mi{place}$ of $o_2$ is also the current activity of
$o_3$.
The lifecycle heuristic with the thresholds $T_{\mi{int}} = 200$ and
$T_{\max} = 2000$ of~\cite{Muehlmeyer2026} learns within this run.
Its first eviction at $i = 7$ agrees with LRU, since $b_4$ is still empty.
This eviction records $\mi{place}$ as the end activity of a completed order
lifecycle.
Hence, the confidence that $\mi{ship}$ ends an order drops, and the expected
lifetime of $o_1$ rises to $1991$, while $i_1$ keeps $1976$.
The heuristic therefore evicts $i_1$ instead of $o_1$ at $i = 7$, and $o_1$ at
$i = 8$.
Both runs yield the same model.
Note that the recorded end activity stems from the premature eviction of
$o_2$, so a premature eviction also misleads the heuristic.
\end{example}

\begin{table}[t]
\centering
\caption{Models of the unbounded run and of the bounded run of
Table~\ref{tab:run} after $s_9$.
They differ only in the start activity $\mi{pay}$ of orders.}
\label{tab:models}
\begin{tabular}{@{}llll@{}}
\toprule
Type & Arcs & Start & End \\
\midrule
$\mi{order}$ (unbounded) & $\mi{place} \to \mi{pay}$,
  $\mi{pay} \to \mi{ship}$ & $\mi{place}$ & $\mi{place}$, $\mi{ship}$ \\
$\mi{order}$ (bounded) & $\mi{place} \to \mi{pay}$,
  $\mi{pay} \to \mi{ship}$ & $\mi{place}$, $\mi{pay}$ &
  $\mi{place}$, $\mi{ship}$ \\
$\mi{item}$ (both) & $\mi{place} \to \mi{pick}$,
  $\mi{pick} \to \mi{ship}$ & $\mi{place}$ & $\mi{place}$, $\mi{ship}$ \\
\bottomrule
\end{tabular}
\end{table}

\subsection{Correctness}
\label{sec:framework:correctness}

The bounded run should yield the same models as the unbounded run.
Whether the unbounded run yields the offline model is a property of the
algorithm, which holds for $A_{\mi{df}}$ and, up to hash collisions,
for~\cite{Muehlmeyer2026} (Appendix~\ref{app:framework}).
Evictions endanger this if an evicted key is needed again.

\begin{definition}[Premature Eviction]
\label{def:premature}
Let $\la q_0, q_1, \dots \ra$ be the run of a streaming algorithm $A$ on a
stream $S$ under a configuration $C$.
Let $k$ be evicted from $B_j$ while processing $s_i$.
The eviction is premature if and only if
$k \in \mi{used}_j(q_{i'-1}, s_{i'})$ for some $i' > i$.
The run is free of premature evictions up to $i$ if and only if no eviction
while processing $s_1, \dots, s_{i-1}$ is premature because of an item
$s_{i'}$ with $i' \le i$.
\end{definition}

An evicted key is missing, so a premature eviction is one whose key is created
again later.
In Example~\ref{ex:run}, the eviction of $o_2$ at $i = 7$ is premature, while the evictions of $o_1$ and $i_1$ are not.

\begin{proposition}[Correctness of Bounded Runs]
\label{prop:correctness}
Let $A = (B_1, \dots, B_n, \mi{upd}, \mi{fin}, \mi{model})$ be a streaming
algorithm with states $Q$.
Buffer $B_j$ is a sink if $\mi{upd}(q, s)$ and $q$ coincide on $B_j$ for all
$q \in Q$ and $s \in \Uitem$.
Let $S$ be an object-centric event stream and let
$C = ((c_1, h_1), \dots, (c_n, h_n))$ be a configuration that fits $S$ with
$c_j = \infty$ for every sink $B_j$.
Let $\la q_0, q_1, \dots \ra$ be the run of $A$ on $S$ under $C$ and
$\la q_0^\infty, q_1^\infty, \dots \ra$ the unbounded run.
Assume that for all $q = (b_1, \dots, b_n) \in Q$, $s \in \Uitem$,
$1 \le j \le n$, and $k \in \mi{dom}(b_j)$:
\begin{enumerate}
  \item[(i)] if $k \notin \mi{used}_j(q, s)$, then
        $\mi{upd}(q, s) = \mi{put}_j(\mi{upd}(\mi{del}_j(q, k), s), k,
        b_j(k))$,
  \item[(ii)] apart from removing $k$, $\mi{fin}(q, j, k)$ changes only
        sinks, and these changes depend only on $k$, $b_j(k)$, and the sinks,
        and
  \item[(iii)] $\mi{model}(\mi{fin}(q, j, k)) = \mi{model}(q)$ if
        $c_j < \infty$.
\end{enumerate}
If the run is free of premature evictions up to $i$, then
$\mi{model}(q_i) = \mi{model}(q_i^\infty)$.
\end{proposition}

Condition~(i) states that an item neither reads nor changes a key it does not
use.
Since a used key is one the item would re-create, an item that needs an
evicted key creates it again, and a premature eviction cannot go unnoticed.
Condition~(ii) lets evictions record information only where items do not look,
e.g., in $B_4$ of Example~\ref{ex:dfg}.
Condition~(iii) states that evicting a key does not change the model.
Appendix~\ref{app:framework} contains the proof and shows that $A_{\mi{df}}$
satisfies all conditions.

The proposition separates two sources of error.
Buffers that violate~(iii) lose model information with every eviction.
In $A_{\mi{df}}$, these are $B_2$, $B_3$, and $B_4$, whose keys are model
elements.
Their size depends on the numbers of activities and object types, so they can
remain unbounded.
Liss et al.~\cite{Liss2025} instead store the arcs as rows in a bounded buffer,
which violates~(iii).
Buffers that satisfy~(iii), such as $B_1$ or the divergence index
of~\cite{Muehlmeyer2026}, lose information only through premature evictions.
For these, the quality of a heuristic can be measured independently of the
model.
Premature evictions cannot be detected online, but they can be counted
afterwards.
We use the fraction of prematurely evicted keys as a measure in our evaluation.
